% ============================================
% OPTION 3: Clarifying Footnote/Note
% Add this to your document where needed
% ============================================

% Version A: As a footnote (add after first use of either η)
% In main.tex, after equation with labor disutility:

\textbf{Labor disutility}:
\begin{equation}
    v(\ell) = \chi \frac{\ell^{1+\eta}}{1+\eta}
\end{equation}
where $\chi$ is the scale parameter and $\eta = 1/\text{Frisch}$ is the inverse Frisch elasticity.\footnote{Note on notation: We use $\eta$ for the inverse Frisch elasticity (a preference parameter) in this section. In the Appendix (full HANK model), $\eta_{i,t}$ denotes the Lagrange multiplier for the cash-in-advance constraint. These are distinct objects, and the context makes clear which is referenced.}


% ============================================
% Version B: As a separate notation box
% Add this at the beginning of your document or appendix
% ============================================

\begin{center}
\fbox{\begin{minipage}{0.9\textwidth}
\textbf{Notation Clarification:}

This document uses the symbol $\eta$ in two distinct contexts:

\begin{enumerate}
    \item \textbf{Preference parameter} (Sections 2-3): $\eta = 1/\text{Frisch}$ denotes the inverse of the Frisch elasticity of labor supply in the utility function $v(\ell) = \chi \frac{\ell^{1+\eta}}{1+\eta}$.

    \item \textbf{Lagrange multiplier} (Appendix): $\eta_{i,t}$ denotes the household-specific Lagrange multiplier associated with the cash-in-advance constraint in the full HANK model.
\end{enumerate}

These are separate mathematical objects. The preference parameter $\eta$ is a time-invariant scalar, while $\eta_{i,t}$ is a time-varying, household-specific endogenous variable.
\end{minipage}}
\end{center}


% ============================================
% Version C: Inline clarification in the appendix
% Modify the line in hank_model_equations.tex:
% ============================================

where $\eta_{i,t}$ is the Lagrange multiplier associated with the cash-in-advance constraint. \emph{Notation: Here $\eta_{i,t}$ represents a Lagrange multiplier (shadow price), distinct from the preference parameter $\eta$ used for the inverse Frisch elasticity in the labor disutility function elsewhere in this document.}


% ============================================
% Version D: Add to both locations
% ============================================

% In main.tex after the labor disutility equation:
where $\chi$ is the scale parameter and $\eta = 1/\text{Frisch}$ is the inverse Frisch elasticity.$^*$

\vspace{0.2cm}
\noindent $^*$\footnotesize{Not to be confused with $\eta_{i,t}$, the Lagrange multiplier in the Appendix.}

% In hank_model_equations.tex after the KKT conditions:
where $\eta_{i,t}$ is the Lagrange multiplier associated with the cash-in-advance constraint.$^{\dagger}$

\vspace{0.2cm}
\noindent $^{\dagger}$\footnotesize{Not to be confused with $\eta$, the inverse Frisch elasticity parameter.}


% ============================================
% RECOMMENDATION:
% ============================================

% Use Version B (notation box) at the start of Section 2 in main.tex
% This provides upfront clarity without cluttering individual equations
% and makes the distinction explicit for readers.
